% Generated draft: requires mathematical and source review.
\documentclass[11pt]{article}
\usepackage[margin=25mm]{geometry}
\usepackage{fontspec}
\setmainfont{Latin Modern Roman}
\usepackage{amsmath,amssymb,mathtools,unicode-math}
\setmathfont{Latin Modern Math}
\usepackage{enumitem,needspace,xurl,hyperref}
\setlength{\parindent}{0pt}
\setlength{\parskip}{5pt}
\setlength{\emergencystretch}{3em}
\newcommand{\N}{\mathbb{N}}
\newcommand{\Z}{\mathbb{Z}}
\newcommand{\Q}{\mathbb{Q}}
\newcommand{\R}{\mathbb{R}}
\newcommand{\C}{\mathbb{C}}
\newcommand{\F}{\mathbb{F}}
\newcommand{\A}{\mathbb{A}}
\newcommand{\PP}{\mathbb{P}}
\newcommand{\E}{\mathbb{E}}
\newcommand{\eps}{\varepsilon}
\newcommand{\dd}{\,\mathrm{d}}
\begin{document}
\section*{UnsolvedMath ID: 30002542}
\textbf{Source problem code:} catalogue entry
\subsection{Definitions and mathematical statement}
For a positive integer \(n\), write \([n]=\{1,\ldots,n\}\) and \(T_n=\binom{[n]}{3}\). The supplied record refers to compatibility of two members of \(T_n\), described only as being ``defined by simultaneous arithmetic progressions.'' It does not specify the arithmetic-progression condition. Consequently, the intended compatibility relation cannot be recovered from the supplied record.

To express the structure of the extracted assertions, let \(C_n\) denote the unspecified symmetric compatibility relation on distinct members of \(T_n\), and define
\[
M_n(C_n)=\bigl\{\mathcal F\subseteq T_n:\ C_n(A,B)\text{ for all distinct }A,B\in\mathcal F\bigr\}.
\]
Define the maximum cardinality and the families attaining it by
\[
r_n(C_n)=\max_{\mathcal F\in M_n(C_n)}|\mathcal F|,
\qquad
\mathcal L_n(C_n)=\{\mathcal F\in M_n(C_n):|\mathcal F|=r_n(C_n)\}.
\]
Thus ``maximum-size'' means maximum cardinality, rather than merely maximal under inclusion.

The supplied extraction reports the following assertions for the intended relation \(C_n\):
\[
|M_n(C_n)|=2^{\binom{n-1}{2}},
\]
and
\[
|\mathcal L_n(C_n)|=
\begin{cases}
2^{m(m-1)},&n=2m+1,\\
2^{(m-1)(m-2)}(2m-1),&n=2m.
\end{cases}
\]
These are schematic transcriptions, not assertions for arbitrary symmetric relations \(C_n\). The intended ranges of \(n\) and \(m\), including any small-parameter exceptions, are not supplied. The positive-integer convention above defines the notation only; it does not establish the source's intended range of validity.

\textbf{Unresolved source qualification.} The record labels these assertions conjectural and gives status ``open,'' but its statement-verification field says that both assertions are proved later in the same official workshop report and are results rather than open problems. It marks the statement ``unrecoverable'' and promotes no clean replacement. Without the missing compatibility definition and a resolution of this status conflict, this record does not specify a precise open problem.
\subsection{Short English statement}
The extracted claims count all families of mutually compatible triples and those having the largest possible cardinality. The compatibility condition and parameter ranges are missing. Moreover, the supplied verification note says both claims are proved in the source report, conflicting with the supplied conjecture and open-status labels. This is a statement-only transcription requiring human mathematical and source review.
\subsection{Sources}
\begin{itemize}
\item Supplied local problem-registry record, permanent UnsolvedMath ID \(30002542\), exact source code \texttt{catalogue entry}, title \emph{Enumeration of Pairwise-Compatible Triple Families}, in the Oberwolfach Reports Open Problems collection. The record lists proposed year 2014 and names no individual proposer.
\item Supplied report citation: \emph{Enumerative Combinatorics} (2014), Oberwolfach Reports, pp.\ 635--720, DOI \texttt{10.4171/owr/2014/12}. \url{https://doi.org/10.4171/owr/2014/12}. The background describes Oberwolfach Reports as documenting mathematical workshops and their open-problem discussions.
\item Supplied literature source, ``Original Oberwolfach report'': \url{https://ems.press/content/serial-article-files/46502?nt=1}.
\item Supplied literature source, ``Stanley transparency'': \url{https://math.mit.edu/~rstan/transparencies/tidbits2}.
\item Literature assessment dated exactly \texttt{2026-08-21}: the supplied curation labels the assertions open and reports finding no authoritative proof or disproof of the general enumerative and extremal claims. This is supplied dated curation, not independently verified here; the background itself requests independent verification of time-sensitive status and cited claims.
\item The supplied statement-verification field reports checking the official workshop report and finding proofs of both extracted conjectures later in that report. It says the original record remains publicly visible for provenance and no separate clean statement is promoted. This supplied verification claim conflicts with the open-status assessment; it has not been independently verified here. The registry gives update timestamp \texttt{2026-09-16T12:06:53Z}, without separately dating this verification claim.
\end{itemize}
\end{document}
